
%\needspace{4\baselineskip}
\color{uniblue}{\subsection{Контроль синхронизма и напряжений (КСН)}}
\color{black}

\begin{enumerate}[label=\arabic{section}.\arabic{subsection}.\arabic{enumi}, labelsep=4pt, leftmargin=0pt, itemindent=57pt, itemsep=0pt, parsep=0pt] % если что можно крутить itemsep

\item
Функциональный блок контроля синхронизма и напряжений (КСН) предназначен для формирования разрешающих сигналов, соответствующих условиям включения, для АПВ, ПАВ и оперативного включения.

\item
В случае присоединения линии к системе через два выключателя для каждого из них можно задействовать свою логику КСН. Далее приведено описание КСН для выключателя В1, логика КСН для второго выключателя аналогична.

\item
Основные режимы включения и контролируемые для них условия приведены в таблице \ref{ksn:tbl1}. Режимы включения с контролем синхронизма (КС) и с улавливанием синхронизма (УС) вводятся соответствующими программными ключами \textbf{«КС В1»} и \textbf{«УС В1»}.

\item
Если разность частот напряжений шин и линии не превышает уставку ««df макс КС»», то разрешается включение с контролем синхронизма. Для режима включения с КС также контролируется:
\begin{itemize}
    \item наличие напряжения на обоих энергообъектах (шины и линия);
    \item разность напряжений шин и линии, не превышающая по модулю значения заданной уставки;
    \item разность фаз напряжений шин и линии, не превышающая по модулю значения заданной уставки.
\end{itemize}

\item
При выводе программного ключа \textbf{«КС В1»} включение разрешается без контроля частоты, а только по уровню напряжений и разности их фаз.

\item
Если разность частот превышает уставку \textbf{«df макс КС»}, но не выходит за пределы уставки \textbf{«df пред. доп.»}, то разрешается включение с улавливанием синхронизма. Сигнал включения с улавливанием синхронизма формируется с опережением на время включения выключателя. Режим включения с УС можно вывести программным ключом \textbf{«УС В1»}.

\item
Наличие или отсутствие напряжения на линии контролируется измерительными органами максимального и минимального напряжения, пороги срабатывания которых задаются уставками \textbf{«Uнал ЛЭП В1»} и \textbf{«Uотс ЛЭП В1»} соответственно. Пороги срабатывания измерительных органов для контроля напряжения на энергообъекте Э1 (шины) задаются уставками \textbf{«Uнал Э1 В1»} и \textbf{«Uотс Э1 В1»}.

\small
\setlength{\extrarowheight}{0.05cm} % высота ячеек в таблице
\begin{longtable}
{|>{\centering\arraybackslash}m{1.7cm}|>{\raggedright\arraybackslash}m{15.2cm}|}
\caption{Условия включения\hfill\vspace{-0.5\baselineskip}}\label{ksn:tbl1} \\ 
\hhline{|-|-|}
\rowcolor{gray!30}
Режим включения & Условия включения  \\ 
\hhline{|-|-|}
\endfirsthead

\caption*{\hspace{3pt}\emph{Продолжение таблицы \ref{ksn:tbl1}\hfill\vspace{-0.5\baselineskip}}} \\ % сделано по ГОСТ 2.105 п.6.8.7
\hhline{|-|-|}
\rowcolor{gray!30}
Режим включения & Условия включения  \\ 
\hhline{|-|-|}
\endhead
\hhline{|-|-|}

\endfoot
\endlastfoot

ОНЛ+ННЭ1 &
\makecell[lt]{%
– $U_{\textsl{ЛЭП}}$ меньше уставки \textbf{«Uотс ЛЭП В1»}; \\
– $U_{\textsl{Э1}}$ больше уставки \textbf{«Uнал Э1 В1»}. } \\
\hline

ОНЭ1+ННЛ &
\makecell[lt]{%
– $U_{\textsl{ЛЭП}}$ больше уставки \textbf{«Uнал ЛЭП В1»}; \\
– $U_{\textsl{Э1}}$ меньше уставки \textbf{«Uотс Э1 В1»}. } \\
\hline

КС(УС) &
\makecell[lt]{%
1) Включение контролем синхронизма: \\
 – $U_{\textsl{ЛЭП}}$ больше уставки \textbf{«Uнал ЛЭП В1»}; \\
 – $U_{\textsl{Э1}}$ больше уставки \textbf{«Uнал Э1 В1»}; \\
 – разность напряжений линии и энергообъекта Э1 не превышает по модулю значения \\ заданной уставки \textbf{«dU макс КС(УС)»}; \\
 – разность фаз напряжений линии и энергообъекта Э1 не превышает по модулю \\ значения заданной уставки \textbf{«dФ макс КС»}; \\
 – разность частот напряжений линии и энергообъекта Э1 не превышает значения \\ заданной уставки \textbf{«df макс КС»}; \\
2) Включение с улавливанием синхронизма: \\
 – $U_{\textsl{ЛЭП}}$ больше уставки \textbf{«Uнал ЛЭП В1»}; \\
 – $U_{\textsl{Э1}}$ больше уставки \textbf{«Uнал Э1 В1»};\\
 – разность напряжений линии и энергообъекта Э1 не превышает по \\ модулю значения заданной уставки \textbf{«dU макс КС(УС)»}; \\
 – разность частот напряжений линии и энергообъекта Э1 превышает \\ значение уставки \textbf{«df макс КС»}; \\
 – разность частот напряжений линии и энергообъекта Э1 не превышает значения \\ заданной уставки \textbf{«df пред. доп.»}; \\
 – за время \textbf{«Т опер. вкл. В1»} с момента подачи сигнала на включение выключателя \\ расхождение углов между векторами напряжений на линии и энергообъекте Э1 не превысит $10^\circ$.} \\
\hline

\end{longtable}
\setlength{\extrarowheight}{0.0cm} % возврат параметра высота ячеек в таблице
\normalsize

\item
Сигнал разрешения АПВ линии \textit{«Разр. АПВл В1»} формируется при выполнении условий включения в режимах ОНЛ+ННЭ1, КС или УС.

\item
Сигнал разрешения АПВ шин \textit{«Разр. АПВш В1»} формируется при выполнении условий включения в режимах ОНЭ1+ННЛ, КС или УС. 

\item
Сигнал разрешения оперативного включения \textit{«Разр. опер. вкл. В1»} может формироваться от внутренней логики КСН или от внешнего сигнала в зависимости от положения программного переключателя \textbf{«Контр. опер. вкл. В1»}.

\item
В положении \textit{«Внутр.»} оперативное включение разрешается при выполнении условий включения в режимах ОНЛ+ННЭ1 (при опробовании линии), ОНЭ1+ННЛ (при опробовании шин), КС или УС (при замыкании в транзит). Также предусмотрена возможность оперативного включения с контролем отсутствия напряжения на шинах или на линии при введенном программном ключе \textbf{«Опер. вкл. В1 с КОН»}. При отсутствии внешнего сигнала \textit{«Опер. вкл. с КС(УС) В1»} оперативное включение может выполняться без контроля напряжений и условий синхронизма.

\item
В положении \textit{«Внеш.»} оперативное включение разрешается при наличии внешнего сигнала «Внеш. сигнал КС(УС) В1».

\item
Сигнал разрешения полуавтоматического включения \textit{«Разр. ПАВ В1»} формируется при выполнении условий включения в режимах КС или УС.

\item
Функциональная схема логики КСН В1 приведена на рисунке \ref{pict_ksn:KSN}. Функциональная схема логики КСН В2 выполняется аналогично схеме логики КСН В1. Уставки КСН В1 представлены в таблице \ref{app_set:ksn}, уставки КСН В2 аналогичны уставкам КСН В1.

\begin{figure}[H]
\centering
\includegraphics[width=0.85\textwidth,height=0.85\textheight,keepaspectratio]{\fbpath/762.3351 КСН/img/KSN.pdf}
\captionof{figure}{Функциональная схема логики КСН В1}\label{pict_ksn:KSN}
\end{figure}

\end{enumerate}